\section{Eisenstein 级数函子}
\markboth{Eisenstein 级数函子}{}

在经典情形，Eisenstein 级数与常数项是在
\[
  G(K)\backslash G(\A)/G(\mathbb O)
  \quad\text{与}\quad
  T(K)N(\A)\backslash G(\A)/G(\mathbb O)
\]
上的函数空间之间起作用的算子，其定义是沿图
\[
  G(K)\backslash G(\A)/G(\mathbb O)\leftarrow B(K)\backslash G(\A)/G(\mathbb O)
  \to T(K)N(\A)\backslash G(\A)/G(\mathbb O)
\]
作拉回—推前。下面我们定义这一手续的几何类比。

\ob{3.1}{blk:3.1}\textbf{紧化与非紧化的几何 Eisenstein 级数。}回忆紧化 Eisenstein 级数函子的定义：
\[
  \overline{\Eis} := \bfP_!\circ\bfQ^* : D(\bBun_T)\to D(\Bun_G).
\]
它的左伴随是「紧化常数项」函子：
\[
  \overline{\CT} := \bfQ_*\circ\bfP^! : D(\Bun_G)\to D(\bBun_T).
\]

复合函子 $\Eis_! := \overline{\Eis}\circ(\fJ_T)_!$ 是通常的（非紧化）Eisenstein 级数函子。它的左伴随 $\CT$ 同构于 $\fJ_T^{\bar*}\circ\overline{\CT}$。

\begin{annote}
$\fJ_T^{\bar*}$ 是原文未加定义的记号。由类型与伴随关系推断，它应读作 $\fJ_T$ 的「重新正规化的 $\ast$-拉回」，即 $(\fJ_T)_!$ 的部分定义左伴随：$\Eis_! = \overline{\Eis}\circ(\fJ_T)_!$ 的左伴随，正是 $(\fJ_T)_!$ 的左伴随与 $\overline{\CT}$ 的复合。译文照录原记号，不代读者改写。\par\vspace{0.15em}{\raggedleft 【来源：原书 p.5 §3.1（A）；译者按（上下文推断，无外部依据）】\par}
\end{annote}

\ob{3.2}{blk:3.2}\textbf{有理 Eisenstein 级数。}记 $\Eis\rt$ 为函子
\[
  \fP\rt_!\circ(\fQ\rt)^* : D(\Bun\rt_T)\to D(\Bun_G),
\]
并记 $\CT\rt$ 为其右伴随。

我们把这两个函子看作本节开头提到的 Eisenstein 算子与常数项算子的几何对应物。

同义反复地，我们有：
\[
  \overline{\Eis}\simeq\Eis\rt\circ\ind_T
  \quad\text{与}\quad
  \CT\rt\simeq\oblv_T\circ\overline{\CT},
\]
\[
  \Eis_!\simeq\Eis\rt\circ(\fJ\rt_T)_!
  \quad\text{与}\quad
  \CT\simeq(\fJ\rt_T)^!\circ\mathcal T_r .
\]

\begin{annote}
上面两行恒等式中的第二式与函子方向不合：$\overline{\CT}$ 取值于 $D(\bBun_T)$，而 $\oblv_T$ 的定义域是 $D(\Bun\rt_T)$；$(\fJ\rt_T)^!$ 亦然，它只能右接取值于 $D(\Bun\rt_T)$ 的函子。按伴随关系，$\overline{\Eis}\simeq\Eis\rt\circ\ind_T$ 取右伴随给出 $\overline{\CT}\simeq\oblv_T\circ\CT\rt$，而 $\Eis_!\simeq\Eis\rt\circ(\fJ\rt_T)_!$ 取左伴随给出 $\CT\simeq(\fJ\rt_T)^!\circ\CT\rt$。故原文应是把 $\CT\rt$ 与 $\overline{\CT}$ 写反，并把 $\CT\rt$ 误排成了未定义的 $\mathcal T_r$。译文照录原文，以便与底本对读；读者按上述两式理解即可。\par\vspace{0.15em}{\raggedleft 【来源：原书 p.5 §3.2（A）；译者按（类型检查与伴随性验算）】\par}
\end{annote}

\ob{3.3}{blk:3.3}\textbf{常数项与 Eisenstein 级数。}在分析 $D(\Bun_G)$ 时，一个基本的角色由「计算 Eisenstein 级数之间的 Hom」这一能力扮演。由伴随性，我们有
\[
  \Hom_{D(\Bun_G)}\bigl(\Eis_!(\mathcal M_1),\Eis_!(\mathcal M_2)\bigr)
  \simeq \Hom_{D(\Bun_T)}\bigl(\mathcal M_1,\CT\circ\Eis_!(\mathcal M_2)\bigr).
\]
也就是说，我们需要理解单子
\[
  \CT\circ\Eis_! : D(\Bun_T)\to D(\Bun_T)
\]
的结构。

由于函子 $\Eis_!$ 可以分解为 $\Eis\rt\circ(\fJ\rt_T)_!$，我们有单子的典范映射
\[
  \widetilde\Omega\to\CT\circ\Eis_! .
\]

下述命题来自用 Zastava 空间对单子 $\widetilde\Omega$ 的计算，以及引理 2.3.1：

\begin{othm}{命题}{3.3.1}{prop:3.3.1}
函子 $\CT\circ\Eis_!$ 带有按 Weyl 群元素 $w\in W$ 编号的典范滤过，使得：
\begin{enumerate}[label=(\alph*), leftmargin=2.4em, itemsep=0.25em]
\item 首项 $(\CT\circ\Eis_!)_{unit}$ 带有单子结构，且作为单子典范同构于 $\widetilde\Omega$；
\item 末项 $(\CT\circ\Eis_!)_{w_0}$ 典范同构于 $w_0$ 在 $D(\Bun_T)$ 上的作用函子。
\end{enumerate}
\end{othm}

\emph{注。}命题的 (a) 可以这样理解：单子 $\widetilde\Omega$ 是单子 $\CT\circ\Eis_!$ 在关于 $\Bun_T$ 的谱意义上「局部」的最佳逼近。

还要指出：在经典的自守函数论中，单子 $\widetilde\Omega$ 的类比物是 $T$ 所对应的 adelic 商上函数空间的一个自同态，它在某个 Grossen 特征（Größencharakter）上的取值，等于（所有正根之积对应的）阿贝尔 $L$ 函数与其移位变元之比的商。

\begin{annote}
两处需要点破。其一，「按 $w\in W$ 的滤过」是指把 $\CT\circ\Eis_!$ 写成归次为 $w$ 的逐次商之链，首项（归次为 $unit$）是单子 $\widetilde\Omega$，末项（归次为 $w_0$）是 $w_0$ 的作用；这正是 §4.5 在谱侧对 $(\fP\st_N)^!\circ\fP\st_*$ 所证的同型结论。其二，经典类比里的「比值」应读作 $L(\chi,s)/L(\chi,s+?)$ 形状的商：正根之积给出 $\prod_{\alpha>0}(1-\chi(\alpha)t)$ 型的因子，移位变元来自 Weyl 向量 $\rho$。\par\vspace{0.15em}{\raggedleft 【来源：原书 p.5 §3.3（A）；经典自守形式论中 Eisenstein 级数的常数项公式（A）】\par}
\end{annote}
